\section{Conclusion}

DiLu-Ollama makes local-first evaluation of LLM driving agents more scientifically usable by combining simulation, task-conditioned benchmarking, runtime-aware validity handling, and publication-facing analysis. For intelligent-vehicle research, the central message is not only that local language-model policies can be benchmarked, but that the benchmark must also reveal when runtime behavior breaks the comparison itself. In the resulting three-tier highway study, lightweight is the only evidence-ready tier, midclass remains screening-only, and highclass exposes a runtime-scalability boundary rather than a competitive ranking. Within that evidence base, the Phi lightweight pair provides a concrete example that fine-tuning can improve local driving performance without increasing mean latency. The claim is intentionally narrow, but that is the point: fine-tuning gains are only meaningful inside the part of the model space where benchmark validity survives the local runtime regime.
